\uuid{0dEe}
\exo7id{6898}
\titre{exo7 6898}
\auteur{ruette}
\organisation{exo7}
\datecreate{2013-01-24}
\isIndication{false}
\isCorrection{false}
\chapitre{Probabilité discrète}
\sousChapitre{Probabilité conditionnelle}
\cpge{mpsi-mp2i.probabilites-univers-fini.c, pcsi.probabilites-variables-aleatoires.c, ptsi.probabilites-variables-aleatoires.c, bcpst1.concepts-base-probabilites.b, mp-mpi.variables-aleatoires-discretes.c, pc.probabilites-lois-usuelles.c, psi.probabilites-lois-usuelles.c, pt.probabilites-lois-usuelles.d}
\module{Probabilité et statistique}
\niveau{L2}
\difficulte{}

\contenu{
\texte{
Les ampoules de la marque X sont fabriquées dans deux usines, A et
B. 20$\%$ des ampoules de l'usine A et 5$\%$ de l'usine B sont
défectueuses. Chaque semaine l'usine A produit $2n$ ampoules et
l'usine B produit $n$ ampoules (où $n\ge 1$ est un entier).
On tire une une ampoule au hasard dans la production d'une semaine.
}
\begin{enumerate}
    \item \question{Quelle est la probabilité que l'ampoule tirée ne soit pas 
défectueuse~?}
    \item \question{Si l'ampoule tirée est défectueuse, quelle est la probabilité 
 qu'elle vienne de l'usine A~?}
\end{enumerate}
}
